% Independent conventional research fragment; excluded from paper.tex.
% Not Lean-certified. Evidence and priority boundaries: PAPER_CONTEXT.md.
\section{Real transport as a weighted matching problem}
\label{sec:real-matching-transport}

The real four-dimensional formula determines the global minimum in every
finite real dimension for a class of spectral penalties. The reduction
uses the stationary operator's spectral blocks and a real orthogonal
normal form for a complex Hermitian projection.

Let $\rho:[0,2]\to\mathbb R$ be continuously differentiable, with
$\rho(0)=\rho'(0)=0$ and $\rho'(s)>0$ for $s>0$. Set
\[
 \phi(t)=\rho\bigl(2\sin(t/2)\bigr),\qquad 0\leq t\leq\pi.
\]
In dimension $n$, write
$\mathcal F_\rho(W)=\sum_{j=1}^n\rho(\sigma_j(I-W))$.
Assume that for some $c\in(0,\pi)$ the derivative $\phi'$ is strictly
increasing on $[0,c]$ and strictly decreasing on $[c,\pi]$.
These hypotheses include the Schatten powers $\rho(s)=s^p$, $p>1$,
and the Charbonnier penalty. Define the one-angle and two-angle costs
\[
 s_\rho(\theta)=\min\{2\phi(\theta),\rho(2)\},
\]
\[
 e_\rho(\theta,\psi)=
   2\min_{|\theta-\psi|/2\leq\delta\leq(\theta+\psi)/2}
    \left[\phi\left(\frac{\theta+\psi}{2}+\delta\right)
          +\phi\left(\frac{\theta+\psi}{2}-\delta\right)\right].
\]
A matching $M$ on $\{1,\ldots,k\}$ is a set of disjoint unordered
pairs; its unmatched vertices are denoted by $U(M)$.

\begin{theorem}[The exact real matching formula]
\label{thm:real-matching-minimum}
Let $P,Q$ be equal-rank real orthogonal projections on $\mathbb R^n$,
and let $\theta_1,\ldots,\theta_k$ be their positive principal angles,
including multiplicities and right angles. Then
\begin{equation}
 \min_{\substack{W\in O(n)\\WPW^T=Q}}
   \sum_{j=1}^n\rho(\sigma_j(I-W))
 =\min_M\left\{
    \sum_{\{i,j\}\in M}e_\rho(\theta_i,\theta_j)
     +\sum_{i\in U(M)}s_\rho(\theta_i)\right\}.
 \label{eq:real-matching-minimum}
\end{equation}
An attaining transport acts on one principal plane for each unmatched
angle and couples the two principal planes of each matched pair.
It fixes the remaining common summands.
\end{theorem}

We first give the algebraic decomposition used in the proof.

\begin{lemma}[A Hermitian projection under real orthogonal coordinates]
\label{lem:hermitian-projection-real-blocks}
Every complex Hermitian projection is conjugate by a real orthogonal
matrix to a direct sum of scalar blocks $0,1$ and blocks
\[
 \begin{pmatrix}
   \lambda&-i\sqrt{\lambda(1-\lambda)}\\
   i\sqrt{\lambda(1-\lambda)}&1-\lambda
 \end{pmatrix},\qquad 0<\lambda\leq\tfrac12.
\]
\end{lemma}

\begin{proof}
Write the projection as $A=X+iY$, where $X$ is real symmetric and
$Y$ is real skew-symmetric. The identities $A^2=A$ give
\[
 X^2-Y^2=X,\qquad XY+YX=Y.
\]
For real $v$, positivity and contractivity of $A$ give
$0\leq v^TXv\leq\|v\|^2$, so the eigenvalues of $X$ lie in $[0,1]$.
The second identity makes $Y$ map the $\lambda$-eigenspace of $X$
to the $(1-\lambda)$-eigenspace. On the former space,
\[
 -Y^2=\lambda(1-\lambda)I.
\]
For $0<\lambda<1/2$, an orthonormal basis $v_j$ of the first space
is therefore matched with the orthonormal basis
$Yv_j/\sqrt{\lambda(1-\lambda)}$ of the second. Each pair gives
the displayed block. On the $1/2$-eigenspace, $2Y$ is an orthogonal
complex structure; choose successive orthonormal pairs $v,2Yv$.
On the zero and one eigenspaces, skew-symmetry and $Y^2=0$ force
$Y=0$, giving the scalar blocks.
\end{proof}

\begin{proof}[Proof of \cref{thm:real-matching-minimum}]
The feasible set is nonempty and compact. Choose a global minimizer $W$.
Extend $\phi$ evenly and periodically to the real line. Its extension is
continuously differentiable, including at zero and $\pi$.
Let $K$ be the real skew-symmetric spectral gradient of $W$: on a real
rotation block through $t$, its value is $\phi'(t)$ times the quarter-turn
matrix, and on the eigenvalues $1,-1$ it is zero.
It commutes with $W$.

For every real skew-symmetric $H$, the spectral first-variation formula
\citep{lewis1996derivatives} gives
\[
 \left.\frac{d}{dt}\mathcal F_\rho(W\exp(tH))\right|_{t=0}
       =-\operatorname{tr}(KH).
\]
As in the proof of \cref{eq:schatten-first-variation}, one may use a local
logarithm whose cut avoids the finite spectrum and apply the Hermitian
spectral derivative theorem. The periodic extension's continuous
differentiability covers the eigenvalues $1,-1$.
When $H$ commutes with $P$, the curve is feasible. Take $H$ to be the
block diagonal compression of $K$ to $P\oplus P^\perp$.
Then $-\operatorname{tr}(KH)=\|H\|_F^2$, so stationarity forces
both diagonal compressions to vanish. For $J=2P-I$ this says
\[
 JKJ=-K.
\]
Consequently $-K^2$ commutes with $P,W$, and $Q=WPW^T$.

Consider a positive spectral level $-K^2=g^2I$. The operator
$\mathcal I=K/g$ is an orthogonal complex structure, and $J$ is an
antilinear isometric involution for this complex structure. Its fixed
space is the range of $P$. Choosing a real orthonormal basis of that
fixed space identifies the whole level with its complexification,
with $J$ acting as conjugation. The restriction of $W$ is then a
complex unitary matrix $V$.

Every eigenphase of $V$ solves $\phi'(t)=g$ in $(0,\pi)$.
The single-peak hypothesis gives at most two roots, say $a,b$.
If there are two roots, write
\[
 V=e^{ia}I+(e^{ib}-e^{ia})A
\]
with a complex Hermitian projection $A$; if the roots coalesce,
$V$ is scalar. Apply \cref{lem:hermitian-projection-real-blocks}.
Its real change of coordinates preserves $J$ and therefore $P$.
It decomposes the real level into orthogonal $P,W,Q$-invariant
blocks of dimension two or four, with source ranks one or two.
The scalar-root case gives two-dimensional blocks directly.

On $\ker K$, the only eigenvalues of $W$ are $1,-1$, so $W$ is a
reflection. The real decomposition of two projections, applied to
$P$ and $(I-W)/2$, gives common invariant blocks of dimension at most
two. A one-dimensional block with $W=-1$ cannot occur at a global
minimum: replacing it by $+1$ preserves transport and lowers the cost.
The other one-dimensional blocks have zero cost and zero angle.

On a two-dimensional rank-one block, the cost is at least
$s_\rho(\theta)$. Indeed, its four orthogonal line transporters have
costs $2\phi(\theta)$, $2\phi(\pi-\theta)$, and $\rho(2)$ for
each of the two reflections; the second is no smaller than the first.
On a four-dimensional rank-two block, \cref{thm:real-four-dimensional-minimum}
gives the lower bound
\[
 \min\{e_\rho(\theta,\psi),
                 s_\rho(\theta)+s_\rho(\psi)\}.
\]
To see the equivalence with that theorem's formula, note that
$e_\rho\leq2\phi(\theta)+2\phi(\psi)$ and
$e_\rho\leq2\rho(2)$, the latter by choosing the endpoint phase.
For $\theta\geq\psi$, the remaining choice is exactly
$\rho(2)+2\phi(\psi)$.

Principal-angle multisets add across common reducing blocks of $P,Q$.
Selecting a pair or two singletons in each four-dimensional block
therefore gives a matching whose value is no larger than the cost of
$W$. Zero angles can be discarded: $s_\rho(0)=0$ and
$e_\rho(\theta,0)=2\phi(\theta)\geq s_\rho(\theta)$.
This proves the lower bound in \cref{eq:real-matching-minimum}.

Conversely, choose a cosine--sine decomposition of $P,Q$ and any matching.
Use an attaining common rotation from
\cref{thm:real-four-dimensional-minimum} on each matched pair of principal
planes. On each unmatched plane choose the cheaper canonical rotation
or reflection. Fix every remaining common direction. The resulting
orthogonal sum is feasible and has exactly the matching's value,
proving the reverse inequality and attainment.
\end{proof}

Put $g_{ij}=s_\rho(\theta_i)+s_\rho(\theta_j)-e_\rho(\theta_i,\theta_j)$.
Then \cref{eq:real-matching-minimum} is equivalently
\[
 \sum_i s_\rho(\theta_i)-\max_M\sum_{\{i,j\}\in M}g_{ij}.
\]
Edges of nonpositive weight may be omitted. Thus evaluation of the scalar
costs followed by a maximum-weight matching gives an exact global
optimization procedure. At $p=3/2$ and for the Charbonnier penalty the
edge costs are explicit; other Schatten exponents require only the
one-dimensional scalar balance equation.

The same matching solves nearest commuting-orthogonal correction. Given
$U\in O(n)$ and a prescribed projection $P$, put $Q=UPU^T$. The map
$W\mapsto V=W^TU$ is a bijection from transporters of $P$ onto $Q$
to orthogonal matrices commuting with $P$. Since
$U-V=(I-W^T)U$, its displacement has the same singular values as
$I-W$. Thus the formula computes an orthogonal correction that preserves
the prescribed subspace and has the least spectral penalty among all
such corrections.

\begin{corollary}[The exact real nuclear minimum]
\label{cor:real-nuclear-matching}
The matching formula also holds for $\rho(s)=s$, with
\[
 s_1(\theta)=\min\{4\sin(\theta/2),2\},\qquad
 e_1(\theta,\psi)=4\sin\frac{\theta+\psi}{2}.
\]
\end{corollary}

\begin{proof}
For $s\in[0,2]$ the Charbonnier penalties satisfy
$|\rho_\tau(s)-s|\leq\tau$. Their objectives converge uniformly on
$O(n)$, so their constrained minima converge to the nuclear minimum.
The one-angle and two-angle costs converge uniformly as well.
There are finitely many matchings, so the minimum of their costs
converges to the minimum of the limiting costs. The limits are exactly
the two displayed expressions, by \cref{eq:real-four-nuclear} and its
one-angle counterpart.
\end{proof}

\begin{corollary}[When the real nuclear minimum equals the complex one]
\label{cor:real-nuclear-parity}
The real nuclear transport minimum equals $2\sum_i\sin\theta_i$
exactly when each positive acute principal angle has even multiplicity.
Right-angle multiplicities are unrestricted.
The real canonical rotation minimizes the nuclear cost exactly when
there is at most one positive principal angle and that angle is at most
$\pi/3$.
\end{corollary}

\begin{proof}
For $0<\theta<\pi/2$, both choices in
$s_1(\theta)=\min\{4\sin(\theta/2),2\}$ strictly exceed
$2\sin\theta$. At a right angle the reflection attains equality.
Also
\[
 e_1(\theta,\psi)
 =4\sin\frac{\theta+\psi}{2}
 \geq2\sin\theta+2\sin\psi,
\]
with equality for positive angles exactly when $\theta=\psi$.
Equality in the finite matching sum therefore requires and permits
pairing each positive acute angle with an equal angle. The complex
minimum is the displayed sine sum.

If there are two positive angles, their paired cost is strictly less
than the sum of their canonical costs, since
\[
 4\sin\frac{\theta+\psi}{2}
 <4\sin\frac\theta2+4\sin\frac\psi2.
\]
For a single positive angle, the one-plane canonical choice is optimal
exactly when $4\sin(\theta/2)\leq2$, or $\theta\leq\pi/3$.
\end{proof}

\begin{corollary}[The two largest angles determine canonical optimality]
\label{cor:real-smoothed-canonical-top-two}
For the Charbonnier cost, let $\theta_{(1)}\geq\theta_{(2)}$ be the
two largest principal angles, inserting zeros when necessary.
The canonical rotation is globally optimal over the real field exactly
when
\[
 \cos\theta_{(1)}\geq\tfrac12\eta^{-2},\qquad
 \tan\frac{\theta_{(1)}}2\tan\frac{\theta_{(2)}}2
       \leq\frac{\tau}{\sqrt{\tau^2+4}}.
\]
\end{corollary}

\begin{proof}
The matching formula shows that canonical transport is optimal exactly
when neither a single reflection nor a paired rotation lowers its cost.
By \cref{cor:real-four-smoothed}, these conditions are
\[
 \cos\theta_i\geq\tfrac12\eta^{-2},\qquad
 \frac{\cos((\theta_i+\theta_j)/2)}
      {\cos((\theta_i-\theta_j)/2)}\geq\eta^{-2}
 \quad(i\ne j).
\]
Writing $u_i=\tan(\theta_i/2)$ makes the ratio
$(1-u_iu_j)/(1+u_iu_j)$. Its inequality becomes
\[
 u_iu_j\leq\frac{1-\eta^{-2}}{1+\eta^{-2}}
           =\frac{\tau}{\sqrt{\tau^2+4}}.
\]
All $u_i$ are nonnegative and increase with the angle, so the two largest
angles give the strongest pair condition. The largest angle also gives
the strongest reflection condition.
\end{proof}
